% =========================================================================
\section{Condiciones de contorno}\label{sec:contorno}
% =========================================================================
Al estudiar distintas situaciones se observa que el campo eléctrico sufre una discontinuidad al atravesar una superficie cargada con densidad $\sigma$
(un plano, una cáscara esférica). Interesa saber qué ocurre exactamente en esa frontera. Es de esperar que, en un caso real, ninguna propiedad
electromagnética cambie de forma abrupta, sino continuamente a lo largo de una \emph{capa de transición} de anchura $h$: los campos físicos son continuos y con
derivadas continuas. Pero no podemos esperar conocer los detalles de esa transición, y por eso se sustituye la situación real por una idealizada, en la que
$h\to0$ (\cref{fig:capa}). Como consecuencia, los campos pueden mostrar una discontinuidad al pasar la frontera.

% Figura: capa de transición real frente a la discontinuidad idealizada.
\begin{figure}[ht]
\centering
\pgfplotsset{capaax/.style={tema, width=0.44\textwidth, height=3.8cm, xmin=-1.6, xmax=1.6, ymin=-0.1, ymax=1.2,
   xtick=\empty, ytick=\empty, grid=none, xlabel={$x$}, ylabel={$\vc{E}$}, axis lines=left, clip=false}}
\begin{tikzpicture}
  \begin{axis}[capaax, name=R, title={\small\scshape (a)\; real}]
    \fill[azulfondo] (axis cs:-0.4,-0.1) rectangle (axis cs:0.4,1.2);
    \draw[grispizarra, densely dashed, line width=0.4pt] (axis cs:-0.4,-0.1) -- (axis cs:-0.4,1.2);
    \draw[grispizarra, densely dashed, line width=0.4pt] (axis cs:0.4,-0.1) -- (axis cs:0.4,1.2);
    \addplot[azulprof, line width=1.2pt, domain=-1.6:1.6, samples=100] {0.15+0.85/(1+exp(-9*x))};
    \draw[<->, grispizarra, line width=0.4pt] (axis cs:-0.4,-0.02) -- (axis cs:0.4,-0.02);
    \node[font=\scriptsize, text=grispizarra, anchor=north] at (axis cs:0,-0.04) {$h$};
    \node[font=\scriptsize, text=azulprof, align=center] at (axis cs:0,0.55) {capa de\\transición};
  \end{axis}
  \begin{axis}[capaax, at={(R.east)}, anchor=west, xshift=0.8cm, title={\small\scshape (b)\; ideal ($h\to0$)}]
    \draw[grispizarra, densely dashed, line width=0.4pt] (axis cs:0,-0.1) -- (axis cs:0,1.2);
    \addplot[azulprof, line width=1.2pt, const plot, mark=none] coordinates {(-1.6,0.15) (0,0.15) (0,1.0) (1.6,1.0)};
    \node[font=\scriptsize, text=azulprof, anchor=west, align=left] at (axis cs:0.05,0.55) {superficie de\\[-2pt]discontinuidad};
  \end{axis}
\end{tikzpicture}
\caption{Un campo físico varía de forma continua (y con derivadas continuas) a través de una capa
de transición de anchura $h$ (a); como no se conocen los detalles de esa transición, se sustituye por una
superficie de discontinuidad, el límite $h\to0$ (b).}
\label{fig:capa}
\end{figure}


Lo analizamos primero para un campo vectorial genérico $\vc{F}$ y luego particularizamos a $\vE$ y a $\phi$.

% -------------------------------------------------------------------------
\subsection{Hipótesis de regularidad}\label{sec:hipotesis}
% -------------------------------------------------------------------------
\begin{definicion}[Campo con discontinuidad en una superficie]\label{def:salto}
Sea $S$ una superficie que separa dos regiones $V_1$ y $V_2$, y sea $\un$ la normal a $S$ que apunta del medio~1 al~2. Decimos que el campo
$\vc{F}:\R^{3}\to\R^{3}$ \emph{presenta una discontinuidad de salto en $S$} si $\vc{F}\in C^{1}(V_1)\cap C^{1}(V_2)$ y admite límites laterales finitos al
aproximarse a $S$,
\[
  \vc{F}_1(\vc{r})=\lim_{\vc{r}\to S^{-}}\vc{F},\qquad\vc{F}_2(\vc{r})=\lim_{\vc{r}\to S^{+}}\vc{F},\qquad
  \norm{\vc{F}_1}<\infty,\quad\norm{\vc{F}_2}<\infty ,
\]
para los puntos de $S$.
\end{definicion}
En términos intuitivos: $\vc{F}$ puede ser discontinuo al atravesar $S$, pero no puede divergir en $S$; debe permanecer acotado en un entorno de la superficie
y ser regular a cada lado.

\begin{observacion}[Qué permite esta hipótesis]\label{obs:hipotesis}
En la condición que obtendremos con el teorema de Gauss, la superficie lateral del cilindro tiene área $O(h)$ y, como $\vc{F}=O(1)$, su flujo es $O(h)\to0$. En la que se
obtiene con Stokes, los dos lados verticales del circuito tienen longitud $O(h)$ y por tanto $\int\vc{F}\cdot\dl=O(h)\to0$. La misma hipótesis justifica, entonces, que despreciemos
estas partes de la superficie y del contorno.
\end{observacion}

% -------------------------------------------------------------------------
\subsection{Condiciones de contorno de un campo vectorial}\label{sec:contorno-general}
% -------------------------------------------------------------------------
Sea $\vc{F}$ un campo del que conocemos sus fuentes,
\begin{equation}\label{eq:fuentes}
  \nab\cdot\vc{F}(\vc{r})=b(\vc{r}),\qquad\nab\times\vc{F}(\vc{r})=\vc{c}(\vc{r}).
\end{equation}
Estas ecuaciones no se pueden aplicar directamente a la discontinuidad. Las usaremos en la capa de transición, donde todo es continuo y con derivada continua, y después
tomaremos límites cuando $h\to0$.

\begin{teorema}[Condiciones de contorno generales]\label{teo:contorno}
Bajo las hipótesis de la \cref{def:salto}, con $\nab\cdot\vc{F}=b$ y $\nab\times\vc{F}=\vc{c}$ en la capa de transición, los saltos de las componentes normal y
tangencial de $\vc{F}$ al cruzar $S$ son
\begin{equation}\label{eq:contorno-normal}
  \boxed{\;F_{2n}-F_{1n}=(\vc{F}_2-\vc{F}_1)\cdot\un=\lim_{h\to0}\bigl(h\,\nab\cdot\vc{F}\bigr)=\lim_{h\to0}(h\,b)\;}
\end{equation}
\begin{equation}\label{eq:contorno-tangencial}
  \boxed{\;\un\times(\vc{F}_2-\vc{F}_1)=\lim_{h\to0}h\,\vc{c}\quad\Longleftrightarrow\quad
  \vc{F}_{2t}-\vc{F}_{1t}=\lim_{h\to0}h\,\bigl(\nab\times\vc{F}\bigr)\times\un\;}
\end{equation}
\end{teorema}
\begin{proof}
\emph{Componente normal.} Escogemos un cilindro de altura $h$ construido en la capa de transición, con bases de área $\Delta S$ paralelas a la superficie
(\figpanel{fig:contorno}{a}) y normales $\un_2=\un$ y $\un_1=-\un$. Integrando $\nab\cdot\vc{F}$ en el volumen $V$ del cilindro y usando el teorema de la divergencia,
\[
  \iiint_{V}\nab\cdot\vc{F}\,\dV=\oiint_{S_V}\vc{F}\cdot\dS=\iiint_{V}b\,\dV .
\]
Si el cilindro es suficientemente pequeño como para que $\vc{F}$ sea constante en sus bases y $b$ constante en todo el volumen,
\[
  \oiint_{S_V}\vc{F}\cdot\dS=\vc{F}_2\cdot\un\,\Delta S+\vc{F}_1\cdot(-\un)\,\Delta S+\iint_{S_{\mathrm{lat}}}\vc{F}\cdot\dS=b\,\Delta S\,h .
\]
Tomando límites con $h\to0$ la integral lateral se anula (\cref{obs:hipotesis}); sin embargo, como $b$ puede tener un valor muy alto en la discontinuidad
(una singularidad concentrada en la superficie), el producto $b\,\Delta S\,h$ puede ser distinto de cero. Dividiendo por $\Delta S$ se obtiene \cref{eq:contorno-normal}.

\emph{Componente tangencial.} Escogemos un pequeño circuito rectangular $C$ en un plano perpendicular a la superficie (\figpanel{fig:contorno}{b}), con dos lados de longitud $\Delta l$
paralelos a $S$, uno a cada lado de ella, y otros dos de longitud $h$ en la capa de transición. Sea $\un'$ la normal al rectángulo (que sale del papel) y $\vc{t}=\un'\times\un$ el
versor tangente a la superficie en el plano del circuito, de modo que $\Delta\vc{l}_2=\Delta l\,\vc{t}$ y $\Delta\vc{l}_1=-\Delta l\,\vc{t}$. Por el teorema de Stokes,
\[
  \iint_{S_C}(\nab\times\vc{F})\cdot\dS=\oint_{C}\vc{F}\cdot\dl=\iint_{S_C}\vc{c}\cdot\dS,
\]
y suponiendo el rectángulo suficientemente pequeño para tomar $\vc{F}$ y $\vc{c}$ constantes,
\[
  \oint_{C}\vc{F}\cdot\dl=\vc{F}_2\cdot\Delta\vc{l}_2+\vc{F}_1\cdot\Delta\vc{l}_1+\int_{\text{lados verticales}}\vc{F}\cdot\dl=\vc{c}\cdot\un'\,h\,\Delta l .
\]
Los lados verticales miden $O(h)$ y contribuyen $O(h)\to0$. Tomando el límite $h\to0$ queda
$(\vc{F}_2-\vc{F}_1)\cdot\vc{t}\,\Delta l=\lim\vc{c}\cdot\un'\,h\,\Delta l$, es decir,
$(\vc{F}_2-\vc{F}_1)\cdot(\un'\times\un)=\lim\vc{c}\cdot\un'\,h$. Con la propiedad del producto mixto
$\vc{A}\cdot(\vc{B}\times\vc{C})=\vc{B}\cdot(\vc{C}\times\vc{A})$, el primer miembro es
$\un'\cdot\bigl[\un\times(\vc{F}_2-\vc{F}_1)\bigr]$, y sacando factor común a $\un'$,
\[
  \un'\cdot\Bigl[\un\times(\vc{F}_2-\vc{F}_1)-\lim_{h\to0}h\,\vc{c}\Bigr]=0 .
\]
Como el circuito puede orientarse con cualquier $\un'$ tangente a la superficie, se obtiene \cref{eq:contorno-tangencial}. Multiplicando vectorialmente por $\un$ (con
$\un\times[\un\times\vc{u}]=-\vc{u}_t$) se obtiene la forma equivalente para las componentes tangenciales.
\end{proof}

% -------------------------------------------------------------------------
\subsection{Condiciones de contorno del campo y del potencial electrostáticos}\label{sec:contorno-electrico}
% -------------------------------------------------------------------------
\begin{corolario}[Condiciones de contorno de $\vE$]\label{cor:contorno-E}
Sea $\sigma$ la densidad superficial de carga en la frontera, entendida como el límite $\sigma=\lim_{h\to0}(\rho\,h)$ de una densidad volúmica
muy grande en la capa de transición. Entonces
\begin{equation}\label{eq:contorno-E}
  \boxed{\;\un\cdot(\vE_2-\vE_1)=\lim_{h\to0}\Bigl(h\,\frac{\rho}{\eo}\Bigr)=\frac{\sigma}{\eo},\qquad
  \un\times(\vE_2-\vE_1)=\vc{0}\;}
\end{equation}
Es decir: las componentes normales a la superficie de separación son \emph{discontinuas}, con un salto $\sigma/\eo$, si existe densidad superficial de carga en la
frontera, mientras que las componentes tangenciales \emph{siempre se conservan}: $\vE_{2t}=\vE_{1t}$.
\end{corolario}
\begin{proof}
Basta tomar $b=\rho/\eo$ y $\vc{c}=\vc{0}$ (porque $\nab\times\vE=\vc{0}$) en el \cref{teo:contorno}; el límite del segundo miembro de la ecuación del rotacional es $\lim h\cdot0=0$.
\end{proof}

\begin{corolario}[Continuidad del potencial]\label{cor:contorno-phi}
El potencial electrostático es continuo al atravesar cualquier superficie de separación: $\phi_2=\phi_1$.
\end{corolario}
\begin{proof}
Sean $\phi_1$ y $\phi_2$ los valores del potencial a cada lado de la capa de transición. La integral puede hacerse sobre cualquier trayectoria a través de la capa
(\cref{eq:int-linea}); tomamos una perpendicular a la superficie (\figpanel{fig:contorno}{c}):
\[
  \phi_2-\phi_1=-\int_{1}^{2}\vE\cdot\dl=-\int_{1}^{2}E_n\,\dif l=-\langle E_n\rangle\,h,
\]
donde $\langle E_n\rangle$ es el valor medio de la componente normal de $\vE$ en la capa de anchura $h$. Tomando el límite,
$\phi_2-\phi_1=-\lim_{h\to0}\langle E_n\rangle\,h=0$, puesto que el campo eléctrico nunca es infinito (a lo sumo $E_n$ tiene una discontinuidad, pero sigue siendo finito).
\end{proof}

% Figura: construcciones para las condiciones de contorno (cilindro, circuito, camino del potencial).
\begin{figure}[ht]
\centering
\tikzset{medio/.style={font=\small, circle, draw=grispizarra, line width=0.4pt, inner sep=1.2pt, fill=white},
         lamina/.style={fill=grisclaro!60, draw=grispizarra!70, line width=0.4pt}}
% (a) cilindro de altura h
\begin{minipage}[t]{0.98\textwidth}\centering
\begin{tikzpicture}[baseline=(current bounding box.north), x=1cm, y=1cm]
  % vista oblicua
  \fill[lamina] (0,0) -- (3.4,0) -- (4.3,0.9) -- (0.9,0.9) -- cycle;
  \draw[vect, azulprof] (0.55,0.45) -- ++(0,1.5) node[above, font=\small] {$\un$};
  \node[medio] at (0.3,0.12) {2}; \node[medio] at (0.3,-0.35) {1};
  \draw[rojoteo, line width=0.7pt] (2.0,0.2) -- (2.0,0.95); \draw[rojoteo, line width=0.7pt] (2.9,0.2) -- (2.9,0.95);
  \draw[rojoteo, line width=0.7pt, fill=white, fill opacity=0.6] (2.45,0.95) ellipse (0.45 and 0.15);
  \draw[rojoteo, line width=0.7pt] (2.0,0.2) arc (180:360:0.45 and 0.15);
  \draw[rojoteo, line width=0.4pt, densely dashed] (2.0,0.2) arc (180:0:0.45 and 0.15);
  \node[font=\scriptsize, text=rojoteo, anchor=west] at (3.0,0.6) {$h$};
  % corte transversal
  \begin{scope}[shift={(6.3,0.2)}]
    \fill[lamina] (-0.2,-0.04) rectangle (3.6,0.04);
    \draw[grispizarra, line width=0.5pt] (-0.2,0) -- (3.6,0);
    \draw[azulprof, ->, line width=0.8pt] (-0.1,-0.5) -- (-0.1,1.0) node[above, font=\small] {$\un$};
    \node[medio] at (-0.45,0.35) {2}; \node[medio] at (-0.45,-0.35) {1};
    \draw[rojoteo, line width=0.9pt] (1.2,-0.6) rectangle (2.4,0.6);
    \draw[->, line width=0.9pt, azulacero] (1.8,0.6) -- ++(0,0.55) node[right, font=\scriptsize] {$\un_2=\un$};
    \draw[->, line width=0.9pt, azulacero] (1.8,-0.6) -- ++(0,-0.55) node[right, font=\scriptsize] {$\un_1=-\un$};
    \draw[<->, rojoteo, line width=0.5pt] (2.7,-0.6) -- (2.7,0.6) node[midway, right, font=\scriptsize] {$h$};
    \draw[<->, rojoteo, line width=0.5pt] (1.2,-0.85) -- (2.4,-0.85) node[midway, below, font=\scriptsize] {$\Delta S$};
  \end{scope}
\end{tikzpicture}\\[2pt]{\small\scshape (a)\; componente normal: cilindro de altura $h$ en la capa de transición}
\end{minipage}\\[10pt]
% (b) circuito
\begin{minipage}[t]{0.98\textwidth}\centering
\begin{tikzpicture}[baseline=(current bounding box.north), x=1cm, y=1cm]
  \fill[lamina] (0,0) -- (3.4,0) -- (4.3,0.9) -- (0.9,0.9) -- cycle;
  \draw[azulprof, ->, line width=0.8pt] (0.55,0.45) -- ++(0,1.5) node[above, font=\small] {$\un$};
  \node[medio] at (0.3,0.12) {2}; \node[medio] at (0.3,-0.35) {1};
  \draw[rojoteo, line width=0.8pt] (1.8,-0.3) rectangle (3.3,1.1);
  \begin{scope}[shift={(6.3,0.2)}]
    \fill[lamina] (-0.2,-0.04) rectangle (3.6,0.04);
    \draw[azulprof, ->, line width=0.8pt] (-0.1,-0.5) -- (-0.1,1.0) node[above, font=\small] {$\un$};
    \node[medio] at (-0.45,0.35) {2}; \node[medio] at (-0.45,-0.35) {1};
    \draw[rojoteo, line width=0.9pt, orienta] (1.0,-0.6) -- (1.0,0.6) -- (2.6,0.6) -- (2.6,-0.6) -- cycle;
    \draw[->, azulacero, line width=0.9pt] (1.8,0.6) -- ++(0.7,0) node[above, font=\scriptsize] {$\vc{t}$};
    \draw[->, azulacero, line width=0.9pt] (1.8,-0.6) -- ++(-0.7,0);
    \node[font=\scriptsize, text=azulacero, anchor=north] at (1.1,-0.64) {$\Delta\vc{l}_1=-\Delta l\,\vc{t}$};
    \node[font=\scriptsize, text=azulacero, anchor=south west] at (1.4,0.66) {$\Delta\vc{l}_2=\Delta l\,\vc{t}$};
    \node[font=\scriptsize, text=rojoteo, anchor=west] at (2.65,0) {$h$};
    \draw[azulprof, line width=0.7pt] (1.8,0) circle (0.12);
    \fill[azulprof] (1.8,0) circle (0.9pt);
    \node[font=\scriptsize, text=azulprof, anchor=north west] at (1.85,-0.08) {$\un'$};
  \end{scope}
\end{tikzpicture}\\[2pt]{\small\scshape (b)\; componente tangencial: circuito rectangular $C$ ($\un'$ sale de la figura)}
\end{minipage}\\[10pt]
% (c) camino del potencial
\begin{minipage}[t]{0.98\textwidth}\centering
\begin{tikzpicture}[baseline=(current bounding box.north), x=1cm, y=1cm]
  \fill[lamina] (0,0) -- (3.4,0) -- (4.3,0.9) -- (0.9,0.9) -- cycle;
  \draw[azulprof, ->, line width=0.8pt] (0.55,0.45) -- ++(0,1.5) node[above, font=\small] {$\un$};
  \node[medio] at (0.3,0.12) {2}; \node[medio] at (0.3,-0.35) {1};
  \draw[rojoteo, line width=0.9pt] (2.2,-0.55) -- (2.2,0.55);
  \fill[rojoteo] (2.2,0.55) circle (1.6pt) node[above right, font=\small] {$B$};
  \fill[rojoteo] (2.2,-0.55) circle (1.6pt) node[below right, font=\small] {$A$};
  \begin{scope}[shift={(6.3,0.2)}]
    \fill[lamina] (-0.2,-0.04) rectangle (3.6,0.04);
    \draw[azulprof, ->, line width=0.8pt] (-0.1,-0.5) -- (-0.1,1.0) node[above, font=\small] {$\un$};
    \node[medio] at (-0.45,0.35) {2}; \node[medio] at (-0.45,-0.35) {1};
    \draw[rojoteo, line width=0.9pt, orienta] (1.8,-0.6) -- (1.8,0.6);
    \fill[rojoteo] (1.8,0.6) circle (1.6pt) node[above right, font=\small] {$B$};
    \fill[rojoteo] (1.8,-0.6) circle (1.6pt) node[below right, font=\small] {$A$};
    \draw[<->, rojoteo, line width=0.5pt] (2.5,-0.6) -- (2.5,0.6) node[midway, right, font=\scriptsize] {$h$};
  \end{scope}
\end{tikzpicture}\\[2pt]{\small\scshape (c)\; potencial: camino $A\to B$ perpendicular a la superficie}
\end{minipage}
\caption{Construcciones geométricas para obtener las condiciones de contorno (la región 1 está
debajo de la superficie y la 2 encima; $\un$ apunta de 1 a 2). (a)~Un cilindro de bases de área $\Delta S$ paralelas a la
superficie y altura $h$, contenido en la capa de transición, para integrar la divergencia. (b)~Un
rectángulo, con dos lados de longitud $\Delta l$ paralelos a la superficie y dos de longitud $h$ (los
dos lados verticales), para integrar el rotacional. (c)~Un camino perpendicular a la
superficie, de $A$ a $B$, para el potencial.}
\label{fig:contorno}
\end{figure}


\begin{observacion}[Consecuencias y comprobaciones]\label{obs:contorno-ejemplos}
Las condiciones se pueden reformular con el potencial: como $\vE=-\nab\phi$, la continuidad de $\vE_t$ equivale a la de $\phi$ (que ya hemos demostrado) y el salto de la
derivada normal es $\partial_n\phi|_2-\partial_n\phi|_1=-\sigma/\eo$. Algunos ejemplos que se comprueban en los seminarios:
\begin{itemize}
  \item \emph{Plano infinito} (\cref{sem:plano}): $\vE=\frac{\sigma}{2\eo}\,\mathrm{sgn}(z)\,\uz$, de modo que el salto en $z=0$ es $\frac{\sigma}{2\eo}-\bigl(-\frac{\sigma}{2\eo}\bigr)=\frac{\sigma}{\eo}$.
  \item \emph{Cáscara esférica} (\cref{sem:cascara}): $E=0$ para $r<a$ y $E=\frac{Q}{4\pi\eo a^{2}}=\frac{\sigma}{\eo}$ para $r=a^{+}$, mientras que $\phi$ es continuo.
  \item \emph{Lámina con densidad volúmica} (\cref{sem:lamina}): no hay carga superficial ($\sigma=0$) y por eso $\vE$ es \emph{continuo} en las caras de la lámina.
\end{itemize}
\end{observacion}
